\documentclass[aps,preprint]{revtex4}%
\usepackage{amsfonts}
\usepackage{amsmath}
\usepackage{amssymb}
\usepackage{graphicx}%
\setcounter{MaxMatrixCols}{30}
%TCIDATA{OutputFilter=latex2.dll}
%TCIDATA{Version=5.50.0.2960}
%TCIDATA{CSTFile=revtex4.cst}
%TCIDATA{Created=Thursday, May 15, 2008 09:56:10}
%TCIDATA{LastRevised=Friday, October 09, 2015 14:15:22}
%TCIDATA{<META NAME="GraphicsSave" CONTENT="32">}
%TCIDATA{<META NAME="SaveForMode" CONTENT="1">}
%TCIDATA{BibliographyScheme=Manual}
%TCIDATA{<META NAME="DocumentShell" CONTENT="Articles\SW\REVTeX 4">}
%BeginMSIPreambleData
\providecommand{\U}[1]{\protect\rule{.1in}{.1in}}
%EndMSIPreambleData
\newtheorem{theorem}{Theorem}
\newtheorem{acknowledgement}[theorem]{Acknowledgement}
\newtheorem{algorithm}[theorem]{Algorithm}
\newtheorem{axiom}[theorem]{Axiom}
\newtheorem{claim}[theorem]{Claim}
\newtheorem{conclusion}[theorem]{Conclusion}
\newtheorem{condition}[theorem]{Condition}
\newtheorem{conjecture}[theorem]{Conjecture}
\newtheorem{corollary}[theorem]{Corollary}
\newtheorem{criterion}[theorem]{Criterion}
\newtheorem{definition}[theorem]{Definition}
\newtheorem{example}[theorem]{Example}
\newtheorem{exercise}[theorem]{Exercise}
\newtheorem{lemma}[theorem]{Lemma}
\newtheorem{notation}[theorem]{Notation}
\newtheorem{problem}[theorem]{Problem}
\newtheorem{proposition}[theorem]{Proposition}
\newtheorem{remark}[theorem]{Remark}
\newtheorem{solution}[theorem]{Solution}
\newtheorem{summary}[theorem]{Summary}
\newenvironment{proof}[1][Proof]{\noindent\textbf{#1.} }{\ \rule{0.5em}{0.5em}}
\begin{document}
\preprint{HEP/123-qed}
\title[Short title for running header]{Long title that appears on the title page}
\author{First Author}
\affiliation{My Institution}
\author{Second Author}
\affiliation{My Institution}
\author{Third Author}
\affiliation{Other Institution}
\keywords{one two three}
\pacs{PACS number}

\begin{abstract}
Shell document for REV\TeX{} 4.

\end{abstract}
\volumeyear{year}
\volumenumber{number}
\issuenumber{number}
\eid{identifier}
\date[Date text]{date}
\received[Received text]{date}

\revised[Revised text]{date}

\accepted[Accepted text]{date}

\published[Published text]{date}

\startpage{101}
\endpage{102}
\maketitle
\tableofcontents


%

\begin{equation}
\min_{\boldsymbol{g}\left(  \boldsymbol{x}\right)  =\boldsymbol{c}}f\left(
\boldsymbol{x}\right)  ,
\end{equation}
where $\boldsymbol{x\in}W$ and $\boldsymbol{c\in}V$
\[
\boldsymbol{g:}W\rightarrow V.
\]
The Lagrangian is%
\begin{equation}
\mathcal{L}\left(  \boldsymbol{x,\mu}\right)  =f\left(  \boldsymbol{x}\right)
-\left\langle \boldsymbol{\mu}\left\vert \boldsymbol{g}\left(  \boldsymbol{x}%
\right)  -\boldsymbol{c}\right.  \right\rangle
\end{equation}
and the constrained minimum of the constrained problem is an unconstrained
saddle point of the Lagrangian. Necessary conditions (KKT\ conditions) for the
constrained stationary points are
\begin{align}
\nabla_{x}\mathcal{L}\left(  \boldsymbol{x,\mu}\right)   &  =\left\langle
\nabla_{x}f\left(  \boldsymbol{x}\right)  \right\vert -\left\langle
\boldsymbol{\mu}\left\vert \nabla_{x}\boldsymbol{g}\left(  \boldsymbol{x}%
\right)  \right.  \right\vert =\boldsymbol{0}\\
&  \Rightarrow\nabla_{x}f\left(  \boldsymbol{x}\right)  =\left\langle
\boldsymbol{\mu}\left\vert \nabla_{x}\boldsymbol{g}\left(  \boldsymbol{x}%
\right)  \right.  \right\rangle =\sum_{1\leq j\leq\varkappa}\mu_{j}\nabla
_{x}\boldsymbol{g}\left(  \boldsymbol{x}\right)  _{j}=\sum_{1\leq
j\leq\varkappa}\mu_{j}\nabla_{x}g_{j}\left(  \boldsymbol{x}\right)  \\
&  =\boldsymbol{\mu}\nabla_{x}\boldsymbol{g}\left(  \boldsymbol{x}\right)
\end{align}
and
\begin{equation}
\nabla_{\mu}\mathcal{L}\left(  \boldsymbol{x,\mu}\right)  =\boldsymbol{g}%
\left(  \boldsymbol{x}\right)  -\boldsymbol{c}=\boldsymbol{0.}%
\end{equation}%
\[
\nabla_{x}f\left(  \boldsymbol{x}\right)  \left[  \nabla_{x}\boldsymbol{g}%
\left(  \boldsymbol{x}\right)  \right]  ^{-1}=\boldsymbol{\mu}%
\]
The constraint condition is recovered by the max wrt to variations of the
Lagrange parameter.

The changes $\delta\boldsymbol{x}$ must maintain the constraint thus
\begin{align*}
\boldsymbol{g}\left(  \boldsymbol{x}+\delta\boldsymbol{x}\right)   &
=\boldsymbol{g}\left(  \boldsymbol{x}\right)  +\nabla_{x}\boldsymbol{g}\left(
\boldsymbol{x}\right)  \left\vert \delta\boldsymbol{x}\right\rangle
=\boldsymbol{g}\left(  \boldsymbol{x}\right) \\
&  \Rightarrow\left[  \nabla_{x}\boldsymbol{g}\left(  \boldsymbol{x}\right)
\right]  \left\vert \delta\boldsymbol{x}\right\rangle =\mathbf{0}\\
&  \Leftrightarrow\left\vert \delta\boldsymbol{x}\right\rangle \in\ker\left[
\nabla_{x}\boldsymbol{g}\left(  \boldsymbol{x}\right)  \right]
\end{align*}
The linear approximation $\left[  \nabla_{x}\boldsymbol{g}\left(
\boldsymbol{x}\right)  \right]  $ to $\boldsymbol{g}$ at $\boldsymbol{x}$
maps
\begin{align*}
d\boldsymbol{g}\left(  \boldsymbol{x}\right)   &  \equiv\nabla_{x}%
\boldsymbol{g}\left(  \boldsymbol{x}\right)  :\mathcal{T}_{\boldsymbol{x}%
}\rightarrow V;\quad\delta\boldsymbol{x}\in\mathcal{T}_{\boldsymbol{x}}\\
\left[  \nabla_{x}\boldsymbol{g}\left(  \boldsymbol{x}\right)  \right]
\left\vert \delta\boldsymbol{x}\right\rangle  &  =\delta\boldsymbol{c}%
\quad\delta\boldsymbol{c}\in V.
\end{align*}
At stationary points for $\left\Vert \delta\boldsymbol{x}\right\Vert
\leq\varepsilon$
\begin{align}
\mathcal{L}\left(  \boldsymbol{x}+\delta\boldsymbol{x,\mu}\right)   &
=f\left(  \boldsymbol{x}+\delta\boldsymbol{x}\right)  -\left\langle
\boldsymbol{\mu}\left\vert \boldsymbol{g}\left(  \boldsymbol{x}+\delta
\boldsymbol{x}\right)  -\boldsymbol{c}\right.  \right\rangle \nonumber\\
&  =f\left(  \boldsymbol{x}\right)  +\left\langle \left.  \nabla_{x}f\left(
\boldsymbol{x}\right)  \right\vert \delta\boldsymbol{x}\right\rangle
-\left\langle \boldsymbol{\mu}\left\vert \nabla_{x}\boldsymbol{g}\left(
\boldsymbol{x}\right)  \right.  \delta\boldsymbol{x}\right\rangle \nonumber\\
&  =\mathcal{L}\left(  \boldsymbol{x,\mu}\right)
\end{align}
giving%
\begin{equation}
\left\langle \left.  \nabla_{x}f\left(  \boldsymbol{x}\right)  \right\vert
\delta\boldsymbol{x}\right\rangle -\left\langle \boldsymbol{\mu}\left\vert
\nabla_{x}\boldsymbol{g}\left(  \boldsymbol{x}\right)  \right.  \delta
\boldsymbol{x}\right\rangle =0\forall\delta\boldsymbol{x\in}\mathcal{T}%
_{\boldsymbol{x}}.
\end{equation}
The normal space $\mathcal{N}_{\boldsymbol{x}}$ of $\boldsymbol{g}\left(
\boldsymbol{x}\right)  =\boldsymbol{c}$ is given by
\[
\mathcal{N}_{\boldsymbol{x}}=\operatorname{linspan}\left\{  \left.  \nabla
_{x}g_{j}\left(  \boldsymbol{x}\right)  \right\vert 1\leq j\leq\varkappa
\right\}
\]
If $\dim\left\{  \mathcal{N}_{\boldsymbol{x}}\right\}  =\varkappa$ then the
constraints are \emph{regular}.

If $\dim\left\{  \operatorname{Ran}\left\{  \left[  \nabla_{x}\boldsymbol{g}%
\left(  \boldsymbol{x}\right)  \right]  \right\}  \right\}  =\varkappa,$ i.e.
the map $\left[  \nabla_{x}\boldsymbol{g}\left(  \boldsymbol{x}\right)
\right]  $ is surjective, \emph{and }$\boldsymbol{g}\left(  \boldsymbol{x}%
\right)  =\boldsymbol{c}$ then this point is regular \cite{Hestenes}. The
tangent space $\mathcal{T}_{\boldsymbol{x}}$ to the level surface
$\boldsymbol{g}\left(  \boldsymbol{x}\right)  =\boldsymbol{c}$ at this point
is the space of vectors
\[
\mathcal{T}_{\boldsymbol{x}}=\operatorname{linspan}\left\{  h\left\vert
\left\langle \left.  \left[  \nabla_{x}\boldsymbol{g}\left(  \boldsymbol{x}%
\right)  \right]  \right\vert h\right\rangle =0\right.  \right\}
=\ker\left\{  \nabla_{x}\boldsymbol{g}\left(  \boldsymbol{x}\right)  \right\}
\]
and
\begin{equation}
\mathcal{T}_{\boldsymbol{x}}=\operatorname{linspan}\left\{  h\left\vert
h=\mathbf{\dot{r}}\left(  0\right)  ;\boldsymbol{g}\left(  \mathbf{r}\left(
t\right)  \right)  =\boldsymbol{c},\left\vert t\right\vert \leq\delta
;\mathbf{r}\left(  0\right)  =\boldsymbol{c}\right.  \right\}  .
\end{equation}
Clearly
\begin{equation}
\mathcal{T}_{\boldsymbol{x}}=\ker\left\{  \nabla_{x}\boldsymbol{g}\left(
\boldsymbol{x}\right)  \right\}  \subset W.
\end{equation}
For regular points $\ker\left\{  \left[  \nabla_{x}\boldsymbol{g}\left(
\boldsymbol{x}\right)  \right]  \right\}  =\mathcal{T}_{\boldsymbol{x}}$, thus
$\dim\left\{  \mathcal{T}_{\boldsymbol{x}}\right\}  =\left(  r-\varkappa
\right)  $ and
\begin{equation}
\mathbb{R}^{r}=\mathcal{T}_{\boldsymbol{x}}\oplus\mathcal{N}_{\boldsymbol{x}}.
\end{equation}
For non regular points $\mathcal{T}_{\boldsymbol{x}}\oplus\mathcal{N}%
_{\boldsymbol{x}}\subset\mathbb{R}^{r}.$

At the minimum $\nabla_{x}f\left(  \boldsymbol{x}_{\#}\right)  \in
\mathcal{N}_{\boldsymbol{x}_{\#}}$ and variations $\delta\boldsymbol{x\in
}\mathcal{T}_{\boldsymbol{x}_{\#}}$ are orthogonal to $\nabla_{x}f\left(
\boldsymbol{x}_{\#}\right)  $ and thus to first order do not change the value
of $f\left(  \boldsymbol{x}_{\#}\right)  $.

\subsection{Examples}

One constraint $\varkappa=1$

1.%
\begin{align*}
g_{1}\left(  \boldsymbol{x}\right)   &  =c_{1}\equiv\sum_{1\leq j\leq r}%
x_{j}=c\equiv\left(  1,\cdots,1\right)  \left(
\begin{array}
[c]{c}%
x_{1}\\
\vdots\\
x_{r}%
\end{array}
\right)  =c\\
\nabla_{x}g_{1}\left(  \boldsymbol{x}\right)   &  =\left(  \frac{\partial
g_{1}}{\partial x_{1}},\cdots,\frac{\partial g_{1}}{\partial x_{r}}\right)
=\left(  1,\cdots,1\right)
\end{align*}%
\begin{align*}
\mathcal{N}_{\boldsymbol{x}}  &  =\left\{  \left.  \eta\left(  1,\cdots
,1\right)  \right\vert \eta\in\mathbb{R}\right\}  \cong\mathbb{R}\\
\mathcal{T}_{\boldsymbol{x}}  &  =\left\{  \left.  \delta\boldsymbol{x}%
\right\vert \delta\boldsymbol{x}\bot\left(  1,\cdots,1\right)  \right\}
\cong\mathbb{R}^{r-1}%
\end{align*}
hence%
\[
\nabla_{x}f\left(  \boldsymbol{x}\right)  =\left\langle \boldsymbol{\mu
}\left\vert \nabla_{x}\boldsymbol{g}\left(  \boldsymbol{x}\right)  \right.
\right\rangle =\mu\left(  1,\cdots,1\right)  .
\]
So at min
\[
\frac{\partial f\left(  \boldsymbol{x}_{\#}\right)  }{\partial x_{1}}%
=\cdots=\frac{\partial f\left(  \boldsymbol{x}_{\#}\right)  }{\partial x_{r}%
}=\mu
\]
Let
\[
0\leq x_{j}\leq1\quad1\leq j\leq r
\]
so one can set
\[
f\left(  \boldsymbol{\theta}\right)  =f\left(  \boldsymbol{x}\left(
\boldsymbol{\theta}\right)  \right)
\]
where%
\[
x_{j}\left(  \boldsymbol{\theta}\right)  =\cos\theta_{j};\quad0\leq\theta
_{j}\leq\frac{\pi}{2}%
\]%
\begin{align*}
g_{1}\left(  \boldsymbol{\theta}\right)   &  =c_{1}\equiv\sum_{1\leq j\leq
r}\cos\theta_{j}=c\\
\nabla_{\theta}g_{1}\left(  \boldsymbol{\theta}\right)   &  =\left(
\frac{\partial g_{1}}{\partial\theta_{1}},\cdots,\frac{\partial g_{1}%
}{\partial\theta_{r}}\right)  =-\left(  \sin\theta_{1},\cdots,\sin\theta
_{r}\right)
\end{align*}%
\[
\frac{\partial f\left(  \boldsymbol{\theta}\right)  }{\partial\theta_{j}}%
=\sum_{k}\frac{\partial f\left(  \boldsymbol{x}\left(  \boldsymbol{\theta
}\right)  \right)  }{\partial x_{k}}\frac{\partial x_{k}}{\partial\theta_{j}%
}=-\sum_{k}\frac{\partial f\left(  \boldsymbol{x}\left(  \boldsymbol{\theta
}\right)  \right)  }{\partial x_{k}}\sin\theta_{k}\delta_{jk}=-\frac{\partial
f\left(  \boldsymbol{x}\left(  \boldsymbol{\theta}\right)  \right)  }{\partial
x_{j}}\sin\theta_{j}%
\]%
\[
\nabla_{\theta}f\left(  \boldsymbol{x}\right)  =\left\langle \boldsymbol{\mu
}\left\vert \nabla_{\theta}\boldsymbol{g}\left(  \boldsymbol{\theta}\right)
\right.  \right\rangle =-\mu\left(  \sin\theta_{1},\cdots,\sin\theta
_{r}\right)  .
\]%
\[
\frac{\partial f\left(  \boldsymbol{\theta}\right)  }{\partial\theta_{j}}%
=-\mu\sin\theta_{j}=-\frac{\partial f\left(  \boldsymbol{x}\left(
\boldsymbol{\theta}\right)  \right)  }{\partial x_{j}}\sin\theta_{j}%
\]
At a constrained min
\begin{align*}
-\mu\sin\theta_{j}+\frac{\partial f\left(  \boldsymbol{x}\left(
\boldsymbol{\theta}\right)  \right)  }{\partial x_{j}}\sin\theta_{j}  &  =0\\
\sin\theta_{j}\left\{  \frac{\partial f\left(  \boldsymbol{x}\left(
\boldsymbol{\theta}\right)  \right)  }{\partial x_{j}}-\mu\right\}   &  =0\\
\sin\theta_{j}  &  =0\text{ and/or }\frac{\partial f\left(  \boldsymbol{x}%
\left(  \boldsymbol{\theta}\right)  \right)  }{\partial x_{j}}=\mu
\end{align*}
and
\[
\sum_{1\leq j\leq r}\cos\theta_{j}=c
\]
note
\[
\sin\theta_{j}=0\Rightarrow x_{j}=\cos\theta_{j}=1;\quad0\leq\theta_{j}%
\leq\frac{\pi}{2}%
\]


If we use the parameterization
\[
g_{1}\left(  \boldsymbol{\theta}\right)  =\sum_{1\leq j\leq r}\cos^{2}%
\theta_{j}=c
\]
then $\sin\theta_{j}$ is replaced by $\sin2\theta_{j}$ and
\[
\sin2\theta_{j}=0\Rightarrow x_{j}=\cos^{2}\theta_{j}=1\text{ \emph{or }%
}0;\quad0\leq\theta_{j}\leq\frac{\pi}{2}.
\]


The eigenvalues of the matrix $F\left(  \mathbb{D}^{2},\boldsymbol{\varphi
}\right)  $ (the diagonal elements of the matrix $\tilde{F}\left(
\mathbb{D}^{2},\boldsymbol{\varphi}\right)  $ ) are
\begin{align}
\varepsilon_{j}  &  =\tilde{F}\left(  \mathbb{D}^{2},\boldsymbol{\varphi
}\right)  _{jj}=\frac{1}{\left\langle \left.  \Psi\right\vert \Psi
\right\rangle }\left\{  \left\langle \left.  \Psi\right\vert b_{j}^{\dagger
}b_{j}H\Psi\right\rangle -\left\langle \left.  b_{j}\Psi\right\vert Hb_{j}%
\Psi\right\rangle \right\} \nonumber\\
&  =N_{j}\left\{  \mathcal{E}_{N}\left(  \mathbb{D}^{2},\boldsymbol{\varphi
}\right)  -\mathcal{E}_{N-1}\left(  \mathbb{D}^{2},\boldsymbol{\varphi
,}\varphi_{j}\right)  \right\}  ,
\end{align}
where the occupation of the GSO's $\left\{  \left\vert \varsigma
_{p}\right\rangle \right\}  $ is
\begin{equation}
N_{j}=\frac{\left\langle \left.  \Psi\right\vert b_{j}^{\dagger}b_{j}%
\Psi\right\rangle }{\left\langle \left.  \Psi\right\vert \Psi\right\rangle }.
\end{equation}
Hence
\begin{equation}
\varepsilon_{j}=N_{j}I\left(  \varphi_{j}\right)  ,
\end{equation}
where $\left\{  I\left(  \varphi_{j}\right)  \right\}  $ are the ionization
energies associated with the GSO's $\left\{  \left\vert \varphi_{j}%
\right\rangle \right\}  $

Thus if $x_{j}=1$ then $\varepsilon_{j}$ can be anything, while if $0\leq
x_{j}<1$ then $\varepsilon_{j}=\mu.$ If $\varepsilon_{j}$ are the Lagrange
parameters form the orbital variation then
\[
N_{j}I\left(  \varphi_{j}\right)  =\mu=const.
\]
This is the work done (changing units) in removing an electron from the
orbital $\varphi_{j}.$

2.
\begin{align*}
f\left(  \boldsymbol{x}\right)   &  =5x_{1}^{2}+4x_{1}x_{2}+x_{2}^{2};\quad
g\left(  \boldsymbol{x}\right)  =3x_{1}+2x_{2}=-5\\
\nabla_{x}f\left(  \boldsymbol{x}\right)   &  =\left(  10x_{1}+4x_{2}%
,4x_{1}+2x_{2}\right)  \quad\nabla_{x}g\left(  \boldsymbol{x}\right)  =\left(
3,2\right)
\end{align*}
so at min
\begin{align*}
\left(  10x_{1}+4x_{2},4x_{1}+2x_{2}\right)   &  =\mu\left(  3,2\right) \\
\left(  5x_{1}+2x_{2},2x_{1}+1x_{2}\right)   &  =\frac{\mu}{2}\left(
3,2\right) \\
2\left(
\begin{array}
[c]{cc}%
5 & 2\\
2 & 1
\end{array}
\right)  \left(
\begin{array}
[c]{c}%
x_{1}\\
x_{2}%
\end{array}
\right)   &  =\mu\left(
\begin{array}
[c]{c}%
3\\
2
\end{array}
\right) \\
\left(
\begin{array}
[c]{c}%
x_{1}\\
x_{2}%
\end{array}
\right)   &  =\frac{\mu}{2}\left(
\begin{array}
[c]{cc}%
5 & 2\\
2 & 1
\end{array}
\right)  ^{-1}\left(
\begin{array}
[c]{c}%
3\\
2
\end{array}
\right)
\end{align*}
\emph{and}
\begin{align*}
3x_{1}+2x_{2}  &  =-5\\
\left(  3,2\right)  \left(
\begin{array}
[c]{c}%
x_{1}\\
x_{2}%
\end{array}
\right)   &  =-5
\end{align*}
the last condition specifies $\mu$.

3.
\begin{align*}
f\left(  \boldsymbol{x}\right)   &  =5x_{1}^{2}+4x_{1}x_{2}+x_{2}^{2};\\
g_{1}\left(  \boldsymbol{x}\right)   &  =3x_{1}^{2}+2x_{2}^{2}=-5\\
g_{2}\left(  \boldsymbol{x}\right)   &  =4x_{1}^{2}+x_{2}^{2}=-1\\
\nabla_{x}f\left(  \boldsymbol{x}\right)   &  =\left(  10x_{1}+4x_{2}%
,4x_{1}+2x_{2}\right) \\
\nabla_{x}g_{1}\left(  \boldsymbol{x}\right)   &  =\left(  6x_{1}%
,4x_{2}\right) \\
\nabla_{x}g_{2}\left(  \boldsymbol{x}\right)   &  =\left(  8x_{1}%
,2x_{2}\right)
\end{align*}
so at min
\begin{align*}
\left(  10x_{1}+4x_{2},4x_{1}+2x_{2}\right)   &  =\mu\left(  3,2\right) \\
\left(  5x_{1}+2x_{2},2x_{1}+1x_{2}\right)   &  =\frac{\mu}{2}\left(
3,2\right) \\
2\left(
\begin{array}
[c]{cc}%
5 & 2\\
2 & 1
\end{array}
\right)  \left(
\begin{array}
[c]{c}%
x_{1}\\
x_{2}%
\end{array}
\right)   &  =\mu\left(
\begin{array}
[c]{c}%
3\\
2
\end{array}
\right) \\
\left(
\begin{array}
[c]{c}%
x_{1}\\
x_{2}%
\end{array}
\right)   &  =\frac{\mu}{2}\left(
\begin{array}
[c]{cc}%
5 & 2\\
2 & 1
\end{array}
\right)  ^{-1}\left(
\begin{array}
[c]{c}%
3\\
2
\end{array}
\right)
\end{align*}
\emph{and}
\begin{align*}
3x_{1}+2x_{2}  &  =-5\\
\left(  3,2\right)  \left(
\begin{array}
[c]{c}%
x_{1}\\
x_{2}%
\end{array}
\right)   &  =-5
\end{align*}
the last condition specifies $\mu$.


\end{document}